\documentclass[10pt,a4paper]{amsart}
\usepackage[margin=19mm]{geometry}
\usepackage{amsmath,amssymb,mathtools}
\usepackage[scr=rsfs,cal=euler]{mathalfa}
\usepackage{xfrac}
\usepackage[unicode,hidelinks,bookmarks=false]{hyperref}
\newcommand{\ie}{i.e.\ }
\newcommand{\cf}{cf.\ }
\newcommand{\resp}{respectively\ }
\newcommand{\Subsection}{\subsection}
\providecommand{\nobreakdash}{\nobreak\hbox{-}}
\setlength{\emergencystretch}{2em}
\multlinegap0pt
\usepackage{bm}
\usepackage{bbm}
\usepackage{dsfont}
\usepackage{xspace}
\usepackage{cancel}
\usepackage{mathtools}
\usepackage{amsthm}
\makeatletter
\@ifundefined{theorem}{\newtheorem{theorem}{Theorem}}{}
\makeatother
\title[On graph automorphisms related to Snort]{On graph automorphisms related to Snort}
\author{Rylo Ashmore, Beth Ann Austin, Alfie M. Davies, Danny Dyer, William Kellough}
\date{Source: arXiv:2506.20669v1 (CC BY 4.0)}
\begin{document}
\maketitle

\noindent\emph{Source and licence.} On graph automorphisms related to Snort. Version arXiv:2506.20669v1 (CC BY 4.0), distributed under the Creative Commons Attribution 4.0 International licence, as recorded in the arXiv licence metadata for this exact version and shown on the abstract page https://arxiv.org/abs/2506.20669v1. Full rights notice: licenses/arxiv-2506.20669-CC-BY-4.0.txt.\par\smallskip

\noindent\textit{Editorial scope.} Counted unit: the Theorem whose source label is \texttt{thm:matchingCharacterization} in the article (source0.tex lines 257--260, source label \texttt{thm:matchingCharacterization}), delivered together with the article's own complete original proof of it (source0.tex lines 262--307, one \texttt{proof} environment of 46 lines). No mathematics is written, repaired or re-derived: statement and proof are the article's own text line for line. Required context retained uncounted: none. Every definition, display and label of those lines is retained unchanged. Omitted from this package and not redistributed: 1--256, 261--261, 308--1769. Nothing inside a retained excerpt is omitted or altered: every retained body is delivered exactly as the article's lines are written, so no omission receipt of any kind accompanies this package. Citations: the retained excerpts carry no citation command at all, so no bibliography entry of the article is transcribed and no reference is redistributed. Reference adaptation: none was needed and none was made; every reference inside the delivered unit resolves inside it, so no unresolved reference or citation is produced. Printed slips kept literal: no printed word, symbol, hypothesis or number of the counted statement or of its proof was altered. Layout and class: the article's own class is \texttt{article} with packages including fontenc, inputenc, babel, cfr-lm, microtype, complexity, authblk, mathtools, amssymb, amsthm, makecell, tikz, ifsym, url; the wrapper is the lane's pinned \texttt{amsart} editorial wrapper at 10pt on A4, which restates the theorem-like environments the retained text needs and adds the article's own macro definitions that the retained text needs (none), with extra wrapper packages (bm, bbm, dsfont, xspace, cancel, mathtools). The delivered numbering is the unit's own. Assets: no retained span contains a figure environment, a graphics-inclusion command, a PDF-page inclusion, a table or a TikZ picture; no image file is copied into this package, the package declares no asset dependency and no image file is redistributed with it. Licence: version arXiv:2506.20669v1 of this article is distributed under the Creative Commons Attribution 4.0 International licence, as recorded in the arXiv licence metadata for this exact version and shown on the abstract page https://arxiv.org/abs/2506.20669v1. A full rights notice is delivered at licenses/arxiv-2506.20669-CC-BY-4.0.txt. Characters: every retained excerpt is pure ASCII, so the delivered engine, XeLaTeX with its default Latin Modern text font, reports no missing character.
\begin{theorem}\label{thm:matchingCharacterization}
    A graph $G$ has an opposition if and only if the complement $\overline{G}$
    has an opposition matching.
\end{theorem}

\begin{proof}
    Let $G$ have an opposition, say $f$. Then in $\overline{G}$ we construct a
    matching $M$ by matching $v$ to $f(v)$. Since $f$ has order two and no
    fixed points, this is well-defined. Since oppositions match vertices to
    non-adjacent vertices, they must be adjacent in the complement, and thus
    are a valid edge of a matching. Since $f$ is an automorphism, we get that
    it is a bijection, and thus $M$ is a perfect matching. Let $v_1v_2$ and
    $v_3v_4$ be two edges of $M$.

    If $v_1v_3$ is an edge in $G$, then so too is $v_2v_4$ because $f$ is an
    automorphism. By the contrapositive, if $v_2v_4$ is an edge in
    $\overline{G}$, then $v_1v_3$ is an edge in $\overline{G}$. Because $f$ is
    an order two automorphism, this is an if and only if argument. The edge
    pair $v_1v_3$ and $v_2v_4$ either both exist, or both do not exist.
    Similarly, the edge pair $v_1v_4$ and $v_2v_3$ either both exist, or both
    do not exist. If no edge pairs exist in $\overline{G}$, then the subgraph
    $H$ induced by $\{v_1,v_2,v_3,v_4\}$ is just $v_1v_2+v_3v_4$. If exactly
    one edge pair exists in $\overline{G}$, then $H \cong C_4$, and if both
    edge pairs exist in $\overline{G}$, then $H \cong K_4$.

    Conversely, suppose $\overline{G}$ has an opposition matching $M$. Define a
    function $f:V(G)\to V(G)$ by mapping $v\in V(G)$ to the matched vertex in
    $M$. We claim $f$ is an opposition in $G$. Let $v_1v_2$ be an edge of $G$,
    and let $f(v_1)f(v_2)=v_3v_4$. Then $v_1v_3$ and $v_2v_4$ are edges in $M$.
    Since $M$ is an opposition matching, the subgraph $H$ induced in
    $\overline{G}$ is isomorphic to either $K_4$, $C_4$, or just the two edges. 
    \begin{itemize}
        \item
            If $H\cong K_4$, then $v_1v_2\not\in E(G)$, and this contradicts
            $v_1v_2\in E(G)$ assumed.
        \item
            If $H\cong C_4$, then $v_1v_2\in E(G)$ implies that the two edges
            not in $H$ are $v_1v_2$ and $v_3v_4$. Thus $v_3v_4$ is an edge of
            $G$.
        \item
            If $H$ is just $v_1v_3 + v_2v_4$, then $v_3v_4$ is not an edge of
            $H$, and so $v_3v_4\in E(G)$.
    \end{itemize}
    It follows that $f$ is an edge-preserving function. Furthermore, $f$ is a
    bijection because it was constructed through a perfect matching. Thus $f$
    is an automorphism. Because $M$ is a perfect matching, the function $f$ has
    order two and has no fixed points. Since $f$ is based on a perfect
    matching, the edge $vf(v)$ always exists in $\overline{G}$, and thus does
    not exist in $G$. Thus the automorphism does not map vertices to adjacent
    vertices. So $f$ is an opposition of $G$.
\end{proof}

\end{document}
